% ch04: 原文 GLDraft_D.tex 行 733–1033；§2.5–2.8
% 只译散文（标题、正文、图题、脚注）；公式、\label、\ref、\eqref、\cite、图文件名保持不变。

\subsection{余迷向膜的对偶}\label{dualco}

为了能把对偶性用于这个问题，我们还需要一个要素：需要理解 $B$-模型（所在空间为 $\M_H(G^\vee,C)$）中 $A$-膜 $\B_\cc$ 与 $\bB_\cc$ 的对偶对象。确切地说，这里指的是 \cite{H} 中称为 $J$ 的复结构下的 $B$-模型；在此复结构下，$\M_H(G^\vee,C)$ 参数化 $C$ 上结构群为 $G^\vee_\C$ 的平坦丛。我们把平坦丛上的联络记作 $\cA=A+\i\phi$，其中 $(A,\phi)$ 是 Hitchin 方程中出现的酉联络与希格斯场。一般的 $B$-膜是 $\M_H(G^\vee,C)$ 上的凝聚层（coherent sheaf），或凝聚层的复形。然而，$A$-膜 $\B_\cc$ 与 $\bB_\cc$ 具有一些额外性质，这些性质意味着它们的对偶 $B$-膜必定相当特殊。

为解释这一点，我们回顾：希格斯丛模空间是超凯勒（hyperkähler）流形，具有复结构 $I,J,K$，它们满足通常的四元数关系，并配有相应的凯勒形式 $\omega_I,\omega_J,\omega_K$ 与复辛形式 $\Omega_I=\omega_J+\i \omega_K$ 等等。几何朗兰兹对偶一般把带有辛结构 $\omega_K$ 的 $\M_H(G,C)$ 的 $A$-模型，映为复结构 $J$ 下 $\M_H(G^\vee,C)$ 的 $B$-模型。当我们不加说明地谈论 $A$-模型或 $B$-模型时，指的就是这两个模型。

在这两个模型中，一般的膜只不过是适当类型的 $A$-膜或 $B$-膜。然而，几何朗兰兹中许多最重要的膜还具有额外性质。例如，支撑在 $\M_H(G^\vee,C)$ 中一点上的膜是类型 $(B,B,B)$ 的膜，也就是说，它在复结构 $I,J$ 与 $K$（或它们的任意线性组合）下都是 $B$-膜。类型 $(B,B,B)$ 的膜的对偶是类型 $(B,A,A)$ 的膜；在希格斯丛模空间的情形，支撑在一点上的秩 1 膜的对偶，是支撑在 Hitchin 纤维化的一条纤维上、带有秩 1 平坦 $\CP$ 丛的膜。这些膜就是 Hecke 本征层，它们是几何朗兰兹纲领研究的中心对象，将在第 \ref{wkb} 节讨论。就当前情形而言，$\B_\cc$ 与 $\bB_\cc$ 是类型 $(A,B,A)$ 的膜：凭借 Kapustin–Orlov 关于余迷向膜（coisotropic brane）的条件，它们是类型 $I$ 与 $K$ 的 $A$-膜；又因为其 $\CP$ 丛的曲率 $\pm\frac{1}{2}\omega_J$ 在复结构 $J$ 下为 $(1,1)$ 型，所以这些丛在复结构 $J$ 下全纯，从而是类型 $J$ 的 $B$-膜。一般地，类型 $(A,B,A)$ 的膜的对偶仍是类型 $(A,B,A)$ 的膜，因此 $\B_\cc$ 与 $\bB_\cc$ 的对偶将是这种类型的膜。类型 $(A,B,A)$ 的膜中最简单的一类，由复结构 $J$ 下某个复拉格朗日子流形（complex Lagrangian submanifold）的结构层给出；用更物理的语言说，这些是支撑在复拉格朗日子流形上、带有平凡 $\CP$ 丛的膜。确实，$\B_\cc$ 与 $\bB_\cc$ 的对偶正是这种类型。这些对偶最早由 Beilinson 与 Drinfeld \cite{BD}（用一种不同的表述）识别出来，其依据是临界 level $k=-h^\vee$ 的共形场论；规范场论的解释则涉及弦论中 D3–NS5 系统与 D3–D5 系统之间的对偶 \cite{GWknots}。\zy{$A$-膜并不一定支撑在拉格朗日子流形上：Kapustin 与 Orlov 指出，支撑在余迷向子流形上、并带有适当 $\CP$ 丛的膜也可以是 $A$-膜，其曲率须满足与辛形式相容的方程。参见 \cite{KO}；在几何朗兰兹语境中的讨论见 \cite{KW}。}

支撑 $\B_\cc$ 的对偶的复拉格朗日子流形，参数化满足“全纯 oper”条件的平坦 $G^\vee_C$ 丛，我们把它记作 $L_\op$。类似地，$\bB_\cc$ 的对偶支撑在复拉格朗日子流形 $L_\bop$ 上，后者参数化满足“反全纯 oper”条件的平坦 $G^\vee_C$ 丛。

全纯 oper 条件可以相当简洁地表述为对丛的全纯类型、即对 $\cA$ 的 $(0,1)$ 部分的整体约束，本文即采用这种表述；它也可以像第 \ref{localoper} 节那样用更局部的方式表述。这两种表述在数学上都是标准的。在四维规范场论中，$S$-对偶把变形 Neumann 边界条件映为变形的“Nahm 极点”边界条件，后者直接给出局部约束 \cite{GWknots}。

定义 $\M_H(G^\vee,C)$ 的复拉格朗日子流形的一种一般方式，是考虑所有具有某个固定全纯类型的平坦 $G^\vee_C$ 丛 $E^\vee\to C$。指定丛的全纯类型，等价于指定 $\cA$ 的 $(0,1)$ 部分。这里以及下面马上要考虑的几个类似情形中，除要求完整联络平坦之外，我们对联络的 $(1,0)$ 部分不加任何约束。在保持 $E^\vee$ 上联络平坦的前提下，我们可以给联络的 $(1,0)$ 部分加上任意的 $\bar\partial_\cA$-闭形式，它代表 $H^0(C,K_C\otimes \ad(E^\vee))$ 中的一个元素。这个空间的维数是 $\M_H(G^\vee,C)$ 维数的一半，所以具有指定全纯类型的平坦 $G^\vee_\C$ 丛构成 $\M_H(G^\vee,C)$ 的一个中间维数子流形 $L$。$L$ 是复拉格朗日子流形，因为复结构 $J$ 下 $\M_H(G^\vee,C)$ 的全纯辛结构 $\Omega_J$ 是 $\cA$ 的 $(1,0)$ 部分与 $(0,1)$ 部分之间的配对，一旦指定 $(0,1)$ 部分便为零。一旦在 $L$ 中选定一个基点，$L$ 就同构于向量空间 $H^0(C,K_C\otimes \ad(E^\vee))$。

我们可以在 $\M_H(G^\vee,C)$ 的同一复结构下定义第二族复拉格朗日子流形：指定平坦 $G^\vee_\C$ 丛的反全纯结构。这相当于指定 $\cA$ 的 $(1,0)$ 部分而让 $(0,1)$ 部分变化；出于与前面相同的理由，它给出一个复拉格朗日子流形。固定 $\cA$ 的 $(1,0)$ 部分或 $(0,1)$ 部分竟然是复结构 $J$ 下的全纯条件，这也许稍感意外。确实，希格斯丛模空间上的复结构 $J$ 对二维流形 $C$ 的复结构并不敏感，它把 $\cA$ 的 $(1,0)$ 部分与 $(0,1)$ 部分完全对称地对待。

子流形 $L_\op$ 与 $L_\bop$ 可以通过指定平坦 $E^\vee$ 丛的全纯结构或反全纯结构的一种特定选择来定义。先解释 $G_\C=\SL(2,\C)$ 情形的定义。对这个群，一个 oper 是平坦丛 $E^\vee\to C$，它作为全纯丛是 $K_C^{-1/2}$ 被 $K_C^{1/2}$ 的非平凡扩张：
\be\label{nonex}0\to K_C^{1/2}\to E^\vee\to K_C^{-1/2}\to 0. \ee
这样的非平凡扩张在同构意义下唯一。因此，这一全纯类型的平坦丛族构成一个复拉格朗日子流形，我们称它为 $L_\op$。作此定义时，我们选定了 $K_C^{1/2}$，等价地选定了 $C$ 上的一个旋量结构（spin structure）。确实，对 $\SL(2,\C)$ 而言，oper 的定义依赖于这样的旋量结构选择（尽管我们在 $L_\op$ 的记号中不体现这一点）。我们将在第 \ref{toposubt} 节回到这一点。

一旦选定基点，即选定一个具有 oper 类型、带平坦联络 $\cA_0$ 的特定 $\SL(2,\C)$ 丛 $E^\vee$，就可以对 $L_\op$ 给出简单的描述。在把 $E^\vee$ 作为 oper 变形时，我们不妨保持 $\cA_0$ 的 $(0,1)$ 部分不变，因为只要求在复规范变换的意义下保持它不变。但我们可以修改 $\cA_0$ 的 $(1,0)$ 部分。为在保持 $\cA_0$ 平坦的同时作此修改，我们给 $\cA_0$ 加上一个 $\bar\partial_{\cA_0}$-闭的 $(1,0)$-形式，即 $H^0(C,K_C\otimes \ad(E^\vee))$ 中的一个元素。利用正合列 (\ref{nonex})，可以证明 $H^0(C,K_C\otimes \ad(E^\vee))$ 同构于 $C$ 上二次微分构成的空间。这正是 $\SL(2,\C)$ 的 Hitchin 纤维化的底空间，所以 $L_\op$ 同构于 Hitchin 纤维化的底空间。这一同构并非完全典范，因为它依赖于 $L_\op$ 中基点的选取。类似的推理适用于其他群。

类似地，对 $G_\C=\SL(2,\C)$，反全纯 oper（简称反 oper）是平坦丛 $E^\vee$，它作为反全纯丛是 $\bar K_C^{-1/2}$ 被 $\bar K_C^{1/2}$ 的非平凡扩张：
\be\label{onex} 0\to\bar K_C^{1/2}\to E^\vee \to \bar K_C^{-1/2}\to 0. \ee
这类平坦丛构成的族是另一个复拉格朗日子流形，我们称它为 $L_\bop$。它非典范地同构于 Hitchin 纤维化的底空间，但带有相反的复结构。

一般地，若 $G_\C^\vee$ 是单复 Lie 群，则有 Lie 代数“主嵌入”（principal embedding）${\su}(2)\to {\g^\vee}$ 的概念。例如，若 $G^\vee=\SL(n,\C)$，主嵌入使得 $\g^\vee$ 的 $n$ 维表示在 $\su(2)$ 下不可约地变换；相应的主子群是 $\SL(2,\C)$ 或 $\SO(3,\C)$ 的拷贝，取决于 $n$ 为偶数还是奇数。为简洁起见，我们有时忽略这一微妙之处，直接说主 $\SL(2,\C)$ 子群，尽管群的整体形式有时是 $\SO(3,\C)$。

$G_\C^\vee$ 的 oper 是一个平坦 $G_\C^\vee$ 丛 $E_\C^\vee$，它作为全纯丛等价于某个 $\SL(2,\C)$ oper 丛的主嵌入；也就是说，等价于一个秩二丛的主嵌入，而该秩二丛是形如 (\ref{nonex}) 的非平凡扩张。对 $G^\vee_\C=\SL(n,\C)$，这意味着 oper 丛在全纯意义下是这样的非平凡扩张的 $(n-1)$ 次对称张量幂，因而有一个同构于 $K_C^{(n-1)/2}$ 的子丛：
\be\label{nex}0\to K_C^{(n-1)/2}\to E^\vee\to \cdots \ee
（并带有由 $K_C$ 的幂给出的滤过）。同样，$E_\C^\vee$ 上联络的 $(1,0)$ 部分除要求完整联络平坦之外不受限制。类似地，反全纯 $G_\C^\vee$ oper 是一个平坦 $G_\C^\vee$ 丛，它作为反全纯丛等价于某个反全纯 $\SL(2,\C)$ oper 丛的主嵌入。

在 $\sigma$-模型层面，$\B_\cc$ 与 $\bB_\cc$ 的对偶是 $L_\op$ 与 $L_\bop$ 的结构层；也就是说，它们是分别支撑在 $L_\op$ 与 $L_\bop$ 上、带有平凡 $\CP$ 丛的秩 1 膜 $\B_\op$ 与 $\B_\bop$。

\subsection{局部约束}\label{localoper}

为给出 oper 的更局部的刻画，我们先考虑 $G^\vee_\C=\SL(2,\C)$ 的情形。$E^\vee$ 的扩张结构意味着存在 $E^\vee \otimes K_C^{-1/2}$ 的一个整体全纯截面 $s$。把联络的 $(1,0)$ 部分记作 $D$。$\SL(2,\C)$-不变的组合 $s \wedge D s$ 是 $C$ 上的整体全纯函数。这个函数必定非零：若它为零，我们就可以写 $D s= a s$，而 $a$ 将定义 $K_C^{-1/2}$ 上的一个全纯平坦联络，这样的联络并不存在（对亏格大于 1 的 $C$，或在低亏格但带有抛物结构的情形）。

不失一般性，我们可以归一化 $s$，使 $s \wedge D s=1$。这使 $s$ 在相差一个符号的意义下确定。对 $G^\vee_\C=\SL(2,\C)$，全纯 oper 条件的局部版本正是：$E^\vee \otimes K_C^{-1/2}$ 具有一个整体定义的全纯截面，使得 $s \wedge D s=1$。

求导可得 $s \wedge D^2 s =0$，因此 $s$ 满足二阶微分方程
\begin{equation}\label{stresst}
D^2 s + t s=0
\end{equation}
其中 $t$ 是 $C$ 上的某个“经典应力张量”（classical stress tensor）。在 $C$ 上作局部坐标变换时，$t$ 按应力张量而非二次微分的方式变换。绝非巧合的是，方程 (\ref{stresst}) 可以看作二维共形场论中退化场关联函数的 Belavin–Polyakov–Zamolodchikov（BPZ）微分方程的经典极限。\zy{BPZ 方程是 Belavin、Polyakov 与 Zamolodchikov 在二维共形场论中导出的线性微分方程，来自退化场的零矢量条件；方程 (\ref{stresst}) 是它在经典极限下的形式。参见 A. Belavin, A. Polyakov and A. Zamolodchikov, ``Infinite conformal symmetry in two-dimensional quantum field theory,'' \emph{Nucl. Phys.} B241 (1984) 333。} 经典应力张量可用来定义 oper 流形上全纯函数代数的一组生成元，这些函数形如
\be\label{opjuno2} f_{t,\alpha}=\int_C \alpha \,t ,\ee
其中 $\alpha$ 是取值于 $T_C$ 的 $(0,1)$-形式。

\begin{yizhu}[译者按：经典应力张量]
原文特别指出 $t$ 按应力张量而非二次微分的方式变换，这一区别值得说明。二次微分 $\alpha$ 在局部坐标变换 $z\mapsto w$ 下按 $\alpha\mapsto \alpha\,(dz/dw)^2$ 变换，只含雅可比因子；而应力张量的变换律还含有一项与坐标变换的 Schwarz 导数有关的贡献。因此 $t$ 并不整体地是一个二次微分，而应理解为 $C$ 上的射影联络；在已固定 $K_C^{1/2}$ 的前提下，对 $\SL(2,\C)$，oper 正对应于 $C$ 上的射影结构（参见 \cite{BD}）。这也解释了为什么方程 (\ref{stresst}) 具有共形场论中应力张量方程的形式。
\end{yizhu}

$G^\vee_\C=\SL(n,\C)$ 的情形可以类似地分析。此时 $E^\vee$ 的 oper 结构（方程 (\ref{nex})）意味着存在 $E^\vee \otimes K_C^{(1-n)/2}$ 的一个整体全纯截面 $s$。于是 $s \wedge D s \wedge \cdots D^{n-1} s$ 是 $C$ 上的整体全纯函数，它不能为零，理由与我们对 $n=2$ 所解释的类似。我们可以归一化 $s$，使 $s \wedge D s \wedge \cdots D^{n-1} s=1$；这唯一地确定 $s$，至多相差乘以一个 $n$ 次单位根。由于 $0=D(s \wedge D s \wedge \cdots D^{n-1} s)=s\wedge Ds \wedge \cdots \wedge  D^{n-2}s\wedge D^ns$，可知 $s$ 满足 $n$ 阶微分方程
\begin{equation}
D^n s + t_2 D^{n-2} s+  \cdots + t_n s =0.
\end{equation}
我们可以定义 oper 流形上全纯函数代数的一组生成元，这些函数形如
\be\label{opjunon} f_{t_k,\alpha}=\int_C \alpha \,t_k \ee
其中 $\alpha$ 是取值于 $T_C^{k-1}$ 的 $(0,1)$-形式，$k=2,\cdots,n$。

对一般的 $G^\vee_\C$，没有像 $\SL(n,\C)$ 的 $n$ 维表示那样方便的、可供使用的特殊表示。不过，给定 oper 丛 $E^\vee$，我们可以考虑它在 $G^\vee_\C$ 的任一不可约表示 $R$ 中的伴丛 $E^\vee_R$。按 oper 的定义，$E^\vee_R$ 作为全纯丛的结构群约化到一个秩 1 子群 $H_\C\subset G^\vee_\C$；这个子群同构于 $\SL(2,\C)$ 或 $\SO(3,\C)$，取决于 $G^\vee_\C$ 与 $R$。设 $R_n$ 是 $H_\C$ 的 $n$ 维不可约表示（对 $\SL(2,\C)$，$n$ 为任意正整数；对 $\SO(3,\C)$，$n$ 为任意正奇数）。作为 $H_\C$ 的表示，有 $R\cong \oplus_{n=0}^\infty Q_n\otimes R_n$，其中 $Q_n$ 是某些向量空间，除有限个外均为零。实际上，若 $N$ 是使 $Q_n$ 非零的最大整数，则 $Q_N$ 是一维的，我们可以把 $Q_N\otimes R_N$ 换成 $R_N$。于是 $R\cong R_N\oplus_{n=0}^{N-1}Q_n\otimes R_n$。在这个分解中，$R_N$ 关于 $H$ 的 Borel 子群 $B_H$ 的最高权向量，也是 $G^\vee$ 关于 $G^\vee$ 的、包含 $B_H$ 的 Borel 子群 $B_{G^\vee}$ 的最高权向量。伴丛 $E^\vee_R$ 作为全纯丛有类似的分解
\be\label{regasso} E^\vee_{R}= E^\vee_{R,N}\oplus \left(\oplus_{n=1}^{N-1} Q_n\otimes E^\vee_{R_n}\right),  \ee
其中 $E^\vee_{R_n}$ 是与某个 $H_\C$ oper 在 $n$ 维表示下相伴的全纯丛。对每个 $n$，由方程 (\ref{nex}) 可以得到向量空间 $Q_n$ 到 $E^\vee_{R} \otimes K_C^{(1-n)/2}$ 的整体全纯截面空间中的典范像 $s_{R,n}$；等价地，得到全纯映射
\be\label{holmap} s_{R,n}: Q_n \otimes K_C^{(n-1)/2} \to E^\vee_{R}. \ee

这里特别重要的是“最高权”对象 $s_{R,N}$，我们直接把它记作 $s_R$：
\be\label{linmap}s_R:K_C^{(N-1)/2}\to E^\vee_R. \ee
这里 $N$ 的定义是：表示 $R$ 的最高权向量（相对于某个 Borel 子群 $B$）位于主 $\SL(2,\C)$ 子群的一个 $N$ 维表示中，而该子群的 Borel 子群包含于 $B$ 中。对 $G^\vee=\SL(n,\C)$ 及 $R$ 为 $n$ 维表示或其对偶的情形，$N=n$。

与 $s \wedge D s=1$ 条件部分相似的条件是：对 $m< n\leq N$，截面族 $D^m s_{R,n}$ 应在 $C$ 的每一点张成 $E^\vee_{R}$。于是 $D^n s_{R,n}$ 可以用 $m<n$ 的 $D^m s_{R,n}$ 展开，从而产生一组复杂的微分方程（例子见附录 \ref{somex}）。这些微分方程的系数是 oper 流形上的全纯函数，可以表示为一族可观测量 $T_{\cal P}$ 的导数的多项式，而这些可观测量所用的标签，与 $G^\vee_\C$ 的量子 Hitchin 哈密顿量的被积函数所用的标签相同。形如 $D^m s_{R,n} \otimes D^{m'} s_{R',n'}$ 的张量积，也可以在张量积 $R \otimes R'$ 分解中出现的所有 $R''$ 所对应的 $D^{m''} s_{R'',n''}$ 的基上展开。展开式的系数同样是 oper 流形上的全纯函数，并可表示为一族可观测量 $T_{\cal P}$ 的导数的多项式。

这族可观测量 $T_{\cal P}$ 有时称为 $G^\vee_\C$ 的经典 $W$-代数（classical $W$-algebra）。用这种语言说，$s_{R,n}$ 所满足的微分方程是 BPZ 方程的经典类比，而张量积展开则类似于退化场的算子乘积展开（OPE）。\zy{经典 $W$-代数由 Drinfeld 与 Sokolov 在研究 KdV 型方程时引入；本文所说的经典 $W$-代数即由这族可观测量 $T_{\cal P}$ 生成。参见 V. Drinfeld and V. Sokolov, ``Lie algebras and equations of Korteweg–de Vries type,'' \emph{J. Soviet Math.} 30 (1985)。}

在四维规范场论中，变形 Neumann 边界条件与变形的 Nahm 极点边界条件互为 $S$-对偶，后者也是全纯-拓扑（holomorphic–topological）的。这个边界条件涉及规范场在边界处某种指定的奇异性。实际上，物理理论中这些奇异的边界条件，在拓扑扭化后的理论中给出一个 oper 边界条件。场在次主导行为中所含的规范不变信息，由边界局部算符承载；这些算符与 $T_{\cal P}$ 可观测量相匹配，并与变形 Neumann 边界条件处用于定义量子 Hitchin 哈密顿量的相应局部算符 $S$-对偶。\footnote{oper 边界条件处的局部算符还包括 oper 流形上的全纯形式。它们来自 $T_{\cal P}$ 的费米子伴；其对偶是附录 \ref{coho} 中描述的算符 $\cal P'$。} $s_{R,n}$ 也自然地出现在规范场论中，我们将在第 \ref{wilop} 节说明。

最后，为了对 $s_{R,n}$ 所满足的各种关系获得更多直观，我们指出：oper 流形有一个更简单的“经典”对应物，即 Hitchin 纤维化的 Hitchin 截面。对 $\SL(2,\C)$，Hitchin 截面参数化那些 $E$ 为直和 $K_C^{1/2} \oplus K_C^{-1/2}$ 的希格斯丛 $(E,\varphi)$。换言之，若我们在复结构 $I$（而非 $J$）中工作，并要求方程 (\ref{nonex}) 中的扩张是分裂的（而非导致 oper 丛的非分裂扩张），所得到的就是 Hitchin 截面。对任意的 $G^\vee_\C$，Hitchin 截面参数化这样的配对 $(E,\varphi)$：$E$ 由 $K_C^{1/2} \oplus K_C^{-1/2}$ 的主嵌入诱导。对 Hitchin 截面，可以推出与我们已经解释的 oper 情形类似的局部条件，只是用希格斯场 $\varphi$ 代替全纯导数 $D$。例如，希格斯丛模空间的 $\SL(2,\C)$ Hitchin 截面，局部地由 $E$ 的一个全纯截面 $s$ 的存在性刻画，该 $s$ 满足 $s \wedge \varphi s=1$ 以及 $\varphi^2 s = \frac12 \Tr\, \varphi^2 s$（后一方程成立只是因为对 $\SL(2,\C)$ 有 $\varphi^2=\frac12 \Tr\, \varphi^2$）。与 oper 情形比较，$D$ 被 $\varphi$ 代替，应力张量 $t$ 被二次微分 $\frac12\Tr\,\varphi^2$ 代替。

一般地，$G^\vee$ 的表示 $R$ 中的伴丛 $E^\vee_R$ 将具有截面 $s_{R,n}$（$n\leq N$），使得 $m<n$ 时的 $\varphi^m s_{R,n}$ 张成 $E^\vee_{R}$。因此 $\varphi^n s_{R,n}$ 可以用 $m<n$ 的 $\varphi^m s_{R,n}$ 展开；同样，对两个表示 $R$、$R'$，$\varphi^m s_{R,n} \otimes \varphi^{m'} s_{R',n'}$ 可以用 $R \otimes R'$ 中所有 $R''$ 对应的 $\varphi^{m''} s_{R'',n''}$ 展开，其系数由规范不变多项式 $\P(\varphi)$ 构成。oper 的这些关系是它们的一种形变。把这一分析稍加推广，就可以恢复第 \ref{dualco} 节的断言：$L_\op$ 非典范地同构于 Hitchin 纤维化的底空间。

\subsection{若干拓扑上的微妙之处}\label{toposubt}

$G_\C$ oper 的定义依赖于 $C$ 上旋量结构的选择，如果主嵌入的像在 $G_\C^\vee$ 中是 $\SL(2,\C)$ 子群；若像是 $\SO(3,\C)$ 子群，则不依赖。（例如，对 $G_\C^\vee=\SL(n,\C)$，oper 的概念依赖旋量结构的选择，当且仅当 $n$ 为偶数。）对某些群而言，这种对旋量结构的依赖，与“$\B_\op$ 是 $\B_\cc$ 的对偶”这一断言构成一种张力，因为 $\B_\cc$ 的定义似乎并不依赖于旋量结构的选择。这一点的解决，本质上已在 \cite{FG} 第 8 节中描述。我们将对秩 1 的群解释其细节。

在与几何朗兰兹相关的 $\N=4$ 超杨–米尔斯理论扭化版本的语境中，$\SO(3)$ 规范场论的电–磁对偶并不是标准的 $\SU(2)$ 规范场论，而是有时所谓的 $\Spin\cdot\SU(2)$ 规范场论。对任意 $d>0$，群 $\Spin(d)\cdot\SU(2)$ 是 $\SO(d)\times \SO(3)$ 的一个二重覆盖，它限制到任一因子上仍是非平凡的二重覆盖。特别地，$\Spin(d)\cdot\SU(2)$ 有到 $\SO(d)$ 与到 $\SO(3)$ 的投影：
\be\label{dobl} \begin{matrix}& & \Spin(d)\cdot\SU(2) && \cr
                                                 &\swarrow &&\searrow &\cr
                                                  \SO(d) &&&& \SO(3). \end{matrix}\ee
所谓 $d$-流形 $M$ 上的 $\Spin(d)\cdot \SU(2)$ 结构，是指 $M$ 上一个以该群为结构群的主丛，使得到第一个因子 $\SO(d)$ 的投影给出 $M$ 的标架丛（即与 $M$ 的切丛 $TM$ 相伴的主丛）。同样，黎曼流形 $M$ 上的 $\Spin(d)\cdot\SU(2)$ 联络，是以该群为结构群、且限制到第一个因子时等于 $M$ 切丛的 Levi-Civita 联络的联络。

假设 $M$ 具有旋量结构，则 $M$ 上 $\Spin(d)\cdot\SU(2)$ 结构的一个通俗描述如下：一旦在 $M$ 上选定一个旋量结构，$\Spin(d)\cdot   \SU(2)$ 丛就等价于一个 $\SU(2)$ 丛 $E^\vee\to M$；若 $M$ 的旋量结构被一个满足 $\ell^2$ 平凡的线丛 $\ell$ 扭化，则 $E^\vee$ 被 $E^\vee\otimes \ell$ 替换。有了这种刻画，就很清楚：尽管 oper 的簇 $L_\op$ 在 $\SU(2)$ 规范场论中没有典范定义，它在 $\Spin(4)\cdot \SU(2)$ 规范场论中却有典范定义。确实，注意到 $\ell\cong\ell^{-1}$，若 $E^\vee$ 出现在用某个 $K^{1/2}$ 的选择定义 oper 的正合列中，则 $E^\vee\otimes \ell$ 出现在把 $K^{1/2}$ 换成 $K^{1/2}\otimes \ell$ 的类似正合列中。

与 $\Spin(4)\cdot \SU(2)$ 规范场论相对，$\SU(2)$ 规范场论的对偶是什么？答案 \cite{FG} 是：对偶为 $\SO(3)$ 规范场论，但在路径积分的定义中要额外包含因子 $\Delta=(-1)^{\int_M w_2(M) w_2(E)}$。这里 $w_2(M)$ 是 $TM$ 的第二 Stiefel–Whitney 类，$w_2(E)$ 是 $\SO(3)$ 丛 $E\to M$ 的第二 Stiefel–Whitney 类；对定向流形，$w_2(M)$ 是旋量结构存在的障碍。\zy{原文此处作“Stieffel–Whitney”，系“Stiefel–Whitney”之笔误；译文径改。参见 J. Milnor and J. Stasheff, \emph{Characteristic Classes}, Princeton University Press (1974)。} 若 $B$-模型的描述用 $\SU(2)$ 规范场论，因而 $L_\op$ 的定义需要 $C$ 上的旋量结构，则 $A$-模型的描述用带有额外因子 $\Delta$ 的 $\SO(3)$ 规范场论，而该因子必须使得 $\B_\cc$ 的定义同样依赖于 $C$ 上旋量结构的选择。这确实如此，理由如下。

若 $M=\Sigma\times C$，其中 $\Sigma$ 与 $C$ 是无边的定向二维流形，则 $w_2(M)=0$ 且 $\Delta=1$，所以路径积分中 $\Delta$ 因子没有影响。但假设 $M$ 有边界 $\partial M$，并带有 $\B_\cc$ 边界条件。为定义拓扑不变量 $\int_M w_2(M) w_2(E)$，需要沿 $\partial M$ 对类 $w_2(M) w_2(E)$ 作平凡化。$\B_\cc$ 边界条件是规范场的一类自由边界条件，因此在 $\B_\cc$ 边界条件下，沿 $\partial M$ 对 $w_2(E)$ 没有限制。但我们可以通过平凡化 $w_2(M)$，即通过沿 $\partial M$ 选取一个旋量结构，来沿 $\partial M$ 平凡化 $w_2(M)w_2(E)$。在我们的应用中，$\partial M$ 是黎曼曲面 $C$ 与一个可缩一维流形（带状区域的一条边界）的乘积，所需要的是 $C$ 上的旋量结构。于是，$\SU(2)$ 规范场论的 $A$-模型对偶是一个 $\SO(3)$ 规范场论；在这个理论中，尽管表面上并非如此，$\B_\cc$（或类似地 $\b\B_\cc$）的定义需要 $C$ 上旋量结构的选择。

这些评注对所有这样的群都有类似物：其 oper 的定义需要旋量结构的选择。它们在通常的 $\N=4$ 超杨–米尔斯理论物理中没有类似物，因为它们只有在拓扑扭化之后才起作用。未扭化的 $\N=4$ 超杨–米尔斯理论有费米子场，其定义需要 $M$ 上的旋量结构。当存在旋量结构时，$\Spin(4)\cdot \SU(2)$ 规范场论与 $\SU(2)$ 规范场论之间的差别就消失了。同样，旋量结构的选择使 $\Delta$ 平凡化。

\subsection{规范场论中 oper 边界条件的拓扑方面}\label{opergauge}

当 $G^\vee$ 的中心 $\zZ(G^\vee)$ 非平凡时，把 $\B_\cc$ 的对偶描述为膜 $\B_\op$（以及 $\b\B_\cc$ 的类似陈述）需要稍作修正。在 $\M_H(G^\vee,C)$ 的 $\sigma$-模型中，$\B_\op$ 是 $\B_\cc$ 的对偶，但我们应记得，低能描述中还包含一个 $\zZ(G^\vee)$ 规范场。在标以 $\B_\op$ 的边界上，$\zZ(G^\vee)$ 规范场被平凡化。正如马上要解释的，这一条件保证了：当理论在 $I\times C$ 上量子化时（其中 $I$ 是带有 $\B_\op$ 与 $\B_\bop$ 边界条件的区间），理论沿 $I$ 承载一个离散电荷。

按照 \cite{KW} 第 7 节的逻辑，这一条件对偶于以下事实：在 $A$-模型一侧，$C$ 上的 $G$-丛有特征类 $\zeta \in H^2(C,\pi_1(G))$，并由 $\int_C\zeta$ 分类（对 $G=\SO(3)$，$\zeta$ 就是第二 Stieffel–Whitney 类 $w_2(E)$）。

Nahm 极点边界条件是 $\B_\op$ 的规范场论版本 \cite{GWknots}，它把边界处的规范群从 $G^\vee$ 约化到其中心 $\zZ(G^\vee)$。沿边界对中心的平凡化是一个附加条件。

对 $G^\vee=\SL(n,\C)$，oper 条件或 Nahm 极点边界条件保证第 \ref{localoper} 节引入的对象 $s$ 存在，并在该节由条件 $s\wedge Ds \wedge \cdots \wedge D^{n-1}s=1$ 归一化，至多相差一个 $n$ 次单位根。$n$ 次单位根是 $\zZ(G^\vee)$ 的元素，因此，把“$\zZ(G^\vee)$ 规范不变性沿边界被平凡化”这一事实表达出来的一个方便说法是：边界上配有归一化截面 $s$ 的一个特定选择。

当我们在带状区域 $\R\times I$（再乘以黎曼曲面 $C$）上量子化理论，并带有 oper 与反 oper 边界条件时，我们在带的左右边界处有这样的平凡化 $s_\ell$ 与 $s_r$。我们可以用 $\zZ(G^\vee)$ 的一个元素 $b$ 作整体规范变换。它作用在平凡化对上为 $(s_\ell,s_r)\to (b s_\ell, b s_r)$，因此凡以这种方式相差者，皆应视为等价。然而，$\zZ(G^\vee)$ 作用在等价类 $(s_\ell,s_r)$ 上为 $(s_\ell,s_r)\to (b s_\ell,s_r)$，$b\in \zZ(G^\vee)$，这导致 $\zZ(G^\vee)$ 在物理希尔伯特空间上的一个作用。在 $B$-模型描述中 $\zZ(G^\vee)$ 对 $\H$ 的这一作用，对偶于 $A$-模型一侧 $\H$ 按 $\int_C \zeta$ 分次这一事实。

对任意 $G^\vee$，边界上 $\zZ(G^\vee)$ 规范场的平凡化可以用第 \ref{localoper} 节引入的对象 $s_{R,n}$ 来表达，但这不如 $\SL(n,\C)$ 情形那么简单。
